% TOP_PROBLEM: {"id": 1414, "title": "Barnette's Conjecture", "category_id": 3, "category": {"id": 3, "name": "graph_theory", "display_name": "Graph Theory", "description": "Problems involving graphs, networks, and their properties.", "slug": "graph-theory", "order_index": 3, "created_at": "2026-07-31T15:26:25.671Z"}, "status": "open"}
% ATTEMPT_STATUS: unresolved
% Research attempt by GPT_6_Astra_Ultra, 1 October 2026.
% Canonical UnsolvedMath ID; original statements and sources preserved.
% FIRST_POSTED: 2026-10-01
% Imported from: unsolvedmath_additional_500.tex; UnsolvedMath ID 1414
% !TeX program = lualatex
% Compile twice: lualatex 1414.tex
% Self-contained: no external bibliography, images, or \input files.
% Numeric section labels and visible numbers are original UnsolvedMath IDs.
\documentclass[11pt,a4paper]{article}
\usepackage[margin=24mm,headheight=14pt]{geometry}
\usepackage{fontspec}
\setmainfont{Latin Modern Roman}
\setsansfont{Latin Modern Sans}
\setmonofont{Latin Modern Mono}
\usepackage{amsmath,amssymb,mathtools}
\usepackage{unicode-math}
\setmathfont{Latin Modern Math}
\usepackage{microtype}
\usepackage{enumitem,xurl,needspace}
\usepackage[dvipsnames]{xcolor}
\usepackage{hyperref}
\hypersetup{unicode=true,colorlinks=true,linkcolor=MidnightBlue,urlcolor=MidnightBlue,
 pdftitle={UnsolvedMath Problem 1414},pdfauthor={GPT 6 Astra Ultra},bookmarksnumbered=true}
\usepackage{bookmark,tocloft,titlesec,fancyhdr}
\setlength{\cftsecnumwidth}{4.2em}
\cftsetpnumwidth{2.5em}
\cftsetrmarg{4em}
\titleformat{\section}{\Large\bfseries\raggedright}{\thesection}{0.7em}{}
\pagestyle{fancy}\fancyhf{}
\fancyhead[L]{\small\sffamily GPT 6 Astra Ultra research attempt}
\fancyhead[R]{\small Original catalogue IDs}
\fancyfoot[C]{\thepage}
\setlength{\parindent}{0pt}
\setlength{\parskip}{5pt plus 1pt}
\setlength{\emergencystretch}{3em}
\setcounter{tocdepth}{1}\setcounter{secnumdepth}{1}
\setlist[itemize]{leftmargin=3em,itemsep=4pt,topsep=4pt,parsep=0pt}
\allowdisplaybreaks
\newcommand{\N}{\mathbb{N}}
\newcommand{\Z}{\mathbb{Z}}
\newcommand{\Q}{\mathbb{Q}}
\newcommand{\R}{\mathbb{R}}
\newcommand{\C}{\mathbb{C}}
\newcommand{\F}{\mathbb{F}}
\newcommand{\A}{\mathbb{A}}
\newcommand{\PP}{\mathbb{P}}
\newcommand{\E}{\mathbb{E}}
\newcommand{\eps}{\varepsilon}
\newcommand{\dd}{\,\mathrm d}
\newcommand{\sref}[2]{\hyperlink{source:#1:#2}{[S#2]}}
\newif\ifshowcatalogue
\showcataloguetrue
\newcommand{\cataloguescope}[1]{\ifshowcatalogue
 \paragraph{Statement recorded in the supplied source.}
 #1\fi}
\begin{document}
% UnsolvedMath ID 1414; source catalogue code GRAPH-027

\setcounter{section}{1413}

\section{Barnette’s Hamiltonian-Cycle Conjecture}\label{1414}

{\small\sffamily UnsolvedMath ID 1414 · Graph Theory}

\subsection{Definitions and mathematical statement}

\paragraph{Shared notation.}Unless a section says otherwise, $\N=\{1,2,\ldots\}$,
$\Z$ is the ring of integers, $\Q,\R,\C$ are the rational, real, and complex
fields, $[N]=\{1,\ldots,\lfloor N\rfloor\}$, and $\log$ is the natural
logarithm. For real functions of a parameter tending to infinity,
$f\ll g$ and $f=O(g)$ mean $|f|\leq Cg$ eventually for some fixed $C$;
$f\gg g$ reverses this comparison; $f=o(g)$ means $f/g\to0$;
$f\sim g$ means $f/g\to1$; and $f\asymp g$ means both order bounds.
A subscript on $O$ or $\ll$ specifies allowed constant dependence.
The expression $n^{o(1)}$ in an upper bound means growth smaller than
$n^\epsilon$ for each fixed $\epsilon>0$. Local definitions govern any
exception, finite-field convention, or use of the same symbol for another
object. An asymptotic claim is not automatically an effective algorithm.



A graph is cubic if every vertex has degree three, bipartite if its vertices admit two independent colour classes, and three-connected if removing fewer than three vertices leaves it connected. A Hamiltonian cycle visits every vertex once. The assertion quantifies over finite simple planar graphs satisfying all three structural hypotheses. \sref{1414}{1} \sref{1414}{2}

\paragraph{Statement recorded in the supplied source.}
Does every cubic bipartite three-connected planar graph have a Hamiltonian cycle? \sref{1414}{1}

\paragraph{Residual task reported by the archive.}

The embedded review (2026-08-17) records: Prove Hamiltonicity or construct a non-Hamiltonian Barnette graph. \sref{1414}{1}

\subsection{Short English statement}

Must every three-connected cubic bipartite planar graph contain a cycle through all of its vertices?

\subsection{Research attempt}

\paragraph{Scope and selected approach.}
Every hypothesis in the statement is retained. The route is to find a
perfect matching, study its complementary two-factor, and seek moves
merging its cycles. The primary abstract of Bekos, Kaufmann, and Pfister
[E1] reports an approximation involving a subhamiltonian cycle; this
does not say that all edges of a Hamiltonian cycle belong to the
original graph. No such relaxation is substituted for the target here.

\paragraph{A perfect matching exists, with a short proof.}
Let $X,Y$ be the bipartition of a finite cubic bipartite graph. Counting
edges from each side gives $3|X|=3|Y|$, hence $|X|=|Y|$. For any
$S\subseteq X$, the $3|S|$ edges incident with $S$ end in $N(S)$,
whose vertices receive at most three of them each. Thus
$|N(S)|\ge|S|$. Hall's matching theorem now gives a matching saturating
$X$ and therefore both sides. To see the relevant content of Hall's
condition directly, a maximum matching leaving some $x\in X$
unmatched has an alternating-reachability set $S$ of left vertices;
every reachable right vertex must be matched, and its mate is in $S$.
The unmatched starting vertex then gives $|N(S)|<|S|$, a contradiction.

Remove such a perfect matching $M$. Every vertex has degree two in
$F=G-M$, so $F$ is a spanning disjoint union of cycles. Bipartiteness
makes their lengths even. The problem has therefore become: can one
choose $M$ so that $F$ has just one component? The existence of a
perfect matching alone gives no control over that number.

\paragraph{Planarity forces quadrilateral faces, not a merging move.}
Write $n,m,f$ for the numbers of vertices, edges, and faces in a plane
embedding. Cubicity gives $m=3n/2$, and Euler's formula gives
$f=n/2+2$. Since the graph is three-connected, each face boundary is
a cycle; simplicity and bipartiteness make its length even and at
least four. Counting incidences of edges and faces,
\[
       \sum_{\text{faces }Q}(6-|Q|)=6f-2m=12.
\]
Only faces of length four contribute positively, contributing two
each. Faces of length six contribute zero, and larger faces contribute
negatively. Hence at least six quadrilateral faces exist. This is a
useful source of possible local moves, but the face count does not
specify which of their edges belong to a chosen perfect matching.

\paragraph{A cycle-merging lemma with its necessary hypotheses.}
Suppose a four-cycle has consecutive vertices $a,b,c,d$, with
$ab,cd\in M$ and $bc,da\in F$. Replacing $ab,cd$ in $M$ by
$bc,da$ gives another perfect matching $M'$. Its complement replaces
$bc,da$ by $ab,cd$. If the two removed $F$ edges lie in different
cycles of $F$, removing them produces two paths, and the two added
edges join those paths into a single cycle. All other components are
unchanged, so the number of components falls by one.

Thus a matching minimizing the number of complementary cycles has
no alternating four-cycle of this merging kind. The contrapositive is
a genuine exclusion on a possible minimal obstruction. But if $bc$
and $da$ lie in the same cycle, the switch can split that cycle rather
than merge two. If the four-cycle is not alternating with respect to
$M$, the described exchange is not a matching exchange at all.
Consequently ``there is a quadrilateral face'' does not imply ``there
is a move improving the two-factor''. The missing global argument
would have to connect the planar face structure to this matching state.

\paragraph{An infinite family where a full cycle is explicit.}
Take two copies of the even cycle $C_{2r}$, $r\ge2$, with vertices
$a_0,\ldots,a_{2r-1}$ and $b_0,\ldots,b_{2r-1}$, and join $a_i$
to $b_i$. This prism is cubic, planar, and bipartite, the latter
coloring obtained by reversing the colors on the second cycle. It is
three-connected: if at most two deleted vertices are in one layer,
the intact other layer connects every remaining piece by a rung; if
one is deleted from each layer, the two surviving layer paths are
connected by a surviving rung. There are at least four rungs before
deletions, so one remains in the latter case.

The following cycle contains every vertex exactly once:
\[
 a_0,a_1,\ldots,a_{2r-1},b_{2r-1},b_{2r-2},\ldots,b_0,a_0.
\]
Its two cross-layer steps are rungs and its other steps are cycle
edges. This proves the conjecture for all even prisms, rather than
only a few sampled sizes. It also illustrates why a convenient
two-factor consisting of the two layer cycles need not be the best
two-factor: a different complementary matching merges them.

\paragraph{Executed verification and limits.}
The graph script constructs prisms with cycle lengths $4,6,8$. It
checks the displayed Hamiltonian cycles, cubic degrees, and connectivity
after every deletion of zero, one, or two vertices. Their orders are
$8,12,16$. The files are \texttt{checks\_graphs.py} and
\texttt{results\_graphs.json} in
\path{_Local/Erdos_problems/UnsolvedMath/attempt_audit_2026_10_01/number_theory/}.
Planarity and bipartiteness of the whole prism family follow from the
construction above, not from an unreported planarity computation.

\paragraph{Remaining gap and outcome.}
The argument provides a spanning even two-factor in every graph under
consideration, a sufficient local condition for reducing its number of
cycles, and a complete infinite subclass. It does not prove that a
minimal two-factor with several cycles always admits that move or a
sequence of more general moves. Nor does it exhibit a non-Hamiltonian
graph satisfying all four hypotheses. The precise missing step is a
global cycle-merging theorem compatible with planarity, bipartiteness,
and three-connectivity. \textbf{Outcome: UNRESOLVED.}

\subsection{Sources}

{\small

\begin{itemize}

\item[\hypertarget{source:1414:1}{S1}] ULAM.AI, \emph{UnsolvedMath}, supplied \texttt{problems.json}, record 1414 (GRAPH-027), statement and background. The embedded literature-review date is 2026-08-17; that is the archive’s review date, not a fresh certification. Dataset landing page: \url{https://huggingface.co/datasets/ulamai/UnsolvedMath}.

\item[\hypertarget{source:1414:2}{S2}] S. Bekos et al., Approximating Barnette's Conjecture, GD 2025, LIPIcs 357, Article 6. \url{https://drops.dagstuhl.de/entities/document/10.4230/LIPIcs.GD.2025.6}. Bibliographic information imported from the supplied record.

\item[\hypertarget{source:1414:3}{S3}] J. Florek, A sufficient condition for cubic 3-connected plane bipartite graphs to be Hamiltonian, J. Graph Theory 110 (2025), \nolinkurl{doi:10.1002/jgt.23270.} \url{https://arxiv.org/abs/2309.09578}. Bibliographic information imported from the supplied record.


\item[E1] Michael A. Bekos, Michael Kaufmann, and
Maximilian Pfister, \emph{Approximating Barnette's Conjecture}, GD 2025,
LIPIcs 357, Article 6. \url{https://drops.dagstuhl.de/entities/document/10.4230/LIPIcs.GD.2025.6}.
Publisher abstract and author metadata inspected 1 October 2026.
The abstract's subhamiltonian approximation is not a proof of the
Hamiltonian-cycle assertion in the original graph.

\end{itemize}

}

\label{end:1414}
\end{document}
